% Part: sets-functions-relations
% Chapter: functions
% Section: kinds

\documentclass[../../../include/open-logic-section]{subfiles}

\begin{document}

\olfileid{sfr}{fun}{kin}
\olsection{ప్రమేయాల రకాలు}

\begin{explain}
మనకు తరచూ ఎదురయ్యే కొన్ని రకాల ప్రమేయాలకు ఒక వర్గీకరణను
పరిచయం చేయడం ఉపయోగకరం.

మొదట, సహప్రవేశంలోని ప్రతి సభ్యమూ ప్రమేయపు విలువగా వచ్చే
ధర్మం గల ప్రమేయాలను పరిగణిద్దాం. ఇటువంటి ప్రమేయాలను
!!{surjective} అంటాం; వీటిని \olref{fig:surjective}లో చూపినట్లు చిత్రించవచ్చు.

\begin{figure}
  \olasset{assets/diagrams/surjective.tikz}
  \caption{!!^a{surjective} ప్రమేయంలో సహప్రవేశంలోని ప్రతి !!{element}
    ప్రమేయపు విలువగా వస్తుంది.}
  \ollabel{fig:surjective}
\end{figure}
\end{explain}

\begin{defn}[!!^{surjective} ప్రమేయం]
$B$ అనేది $f$ వ్యాప్తి కూడా అయితే, అప్పుడే
$f \colon A \rightarrow B$ అనే ప్రమేయాన్ని \emph{!!{surjective}} అంటాం.
అంటే ప్రతి $y \in B$కూ $f(x) = y$ అయ్యే $x \in A$ కనీసం ఒకటి ఉంటుంది.
సంకేతాలలో:
\[
  (\forall y \in B)(\exists x \in A)f(x) = y.
\]
ఇటువంటి ప్రమేయాన్ని $A$ నుంచి $B$కు !!a{surjection} అంటాం.
\end{defn}

\begin{explain}
$f$ !!a{surjection} అని చూపాలంటే, $f$ సహప్రవేశంలోని ప్రతి వస్తువూ
ఏదో ఒక ఇన్‌పుట్ $x$కు $f(x)$ విలువగా వస్తుందని చూపాలి.

ఏ ప్రమేయమైనా !!a{surjection}ను \emph{ఉత్పన్నం చేస్తుందని} గమనించండి.
$f \colon A \to B$ అనే ప్రమేయం ఇస్తే, $f'(x) = f(x)$ అయ్యేలా
$f' \colon A \to \ran{f}$ను నిర్వచిద్దాం.
$\ran{f}$ను $\Setabs{f(x) \in B}{x \in A}$గా \emph{నిర్వచిస్తాం} కనుక,
ఈ $f'$ ప్రమేయం తప్పక !!a{surjection} అవుతుంది.
\end{explain}

\begin{explain}
ఏ ప్రమేయమైనా సాధ్యమైన ప్రతి ఇన్‌పుట్‌ను ఒకే ఫలితానికి ప్రతిచిత్రిస్తుంది.
అయితే వేర్వేరు ఇన్‌పుట్లను ఎప్పుడూ ఒకే ఫలితానికి ప్రతిచిత్రించని
ప్రమేయాలు కూడా ఉంటాయి. ఇటువంటి ప్రమేయాలను !!{injective} అంటాం;
వీటిని \olref{fig:injective}లో చూపినట్లు చిత్రించవచ్చు.
\begin{figure}
  \olasset{assets/diagrams/injective.tikz}
  \caption{!!^a{injective} ప్రమేయం రెండు వేర్వేరు ఆర్గ్యుమెంట్లను
    ఎప్పుడూ ఒకే విలువకు ప్రతిచిత్రించదు.}
  \ollabel{fig:injective}
\end{figure}
\end{explain}

\begin{defn}[!!^{injective} ప్రమేయం]
ప్రతి $y \in B$కూ, $f(x) = y$ అయ్యే $x \in A$ గరిష్ఠంగా ఒకటే ఉంటే,
అప్పుడే $f \colon A \rightarrow B$ను \emph{!!{injective}} అంటాం.
ఇటువంటి ప్రమేయాన్ని $A$ నుంచి $B$కు !!a{injection} అంటాం.
\end{defn}

\begin{explain}
$f$ !!a{injection} అని చూపాలంటే, $f$ ప్రవేశంలోని ఏ
!!{element}s $x$, $y$కైనా, $f(x)=f(y)$ అయితే $x=y$ అని చూపాలి.
\end{explain}

\begin{ex}
$f(x) = 1$ ద్వారా ఇచ్చే స్థిర ప్రమేయం $f\colon \Nat \to \Nat$,
!!{injective} కాదు, !!{surjective} కూడా కాదు.

$f(x) = x$ ద్వారా ఇచ్చే తాదాత్మ్య ప్రమేయం $f\colon \Nat \to \Nat$,
!!{injective}, !!{surjective} రెండూ అవుతుంది.

$f(x) = x+1$ ద్వారా ఇచ్చే ఉత్తరవర్తి ప్రమేయం $f \colon \Nat \to \Nat$,
!!{injective} అవుతుంది కానీ !!{surjective} కాదు.

కింది విధంగా నిర్వచించిన $f \colon \Nat \to \Nat$ అనే ప్రమేయం:
\[
  f(x) =
  \begin{cases}
    \frac{x}{2} & \text{$x$ సరి సంఖ్య అయితే} \\
    \frac{x+1}{2} & \text{$x$ బేసి సంఖ్య అయితే.}
  \end{cases}
\]
!!{surjective} అవుతుంది కానీ !!{injective} కాదు.
\end{ex}

\begin{explain}
!!{injective}, !!{surjective} రెండూ అయ్యే ప్రమేయాలను మనం తరచూ
పరిగణించాల్సి వస్తుంది. ఇటువంటి ప్రమేయాలను !!{bijective} అంటాం.
ఇవి \olref{fig:bijective}లో చిత్రించిన ప్రమేయంలా ఉంటాయి.
!!^{bijection}sను కొన్నిసార్లు \emph{ఒకటి-ఒకటి అనురూపతలు} అని కూడా
అంటాం: ఇవి సహప్రవేశంలోని మూలకాలను ప్రవేశంలోని మూలకాలతో ఏకైకంగా జతచేస్తాయి.
\begin{figure}
  \olasset{assets/diagrams/bijective.tikz}
  \caption{!!^a{bijective} ప్రమేయం సహప్రవేశంలోని మూలకాలను
    ప్రవేశంలోని మూలకాలతో ఏకైకంగా జతచేస్తుంది.}
  \ollabel{fig:bijective}
\end{figure}
\end{explain}

\begin{defn}[!!^{bijection}]
$f \colon A \to B$ అనే ప్రమేయం !!{surjective}, !!{injective} రెండూ
అయితే, అప్పుడే అది \emph{!!{bijective}} అవుతుంది.
ఇటువంటి ప్రమేయాన్ని $A$ నుంచి $B$కు (లేదా $A$, $B$ల మధ్య)
!!a{bijection} అంటాం.
\end{defn}

\end{document}
